%==============================================================================
%  Homework template  --  999 Controls
%
%  To reuse for the next assignment:
%    1. copy this folder,  2. edit the METADATA block below,
%    3. edit the SOLUTIONS section,  4. run  `make`.
%
%  Build:  make          (regenerates code output, then builds main.pdf)
%          make clean    (removes build artifacts)
%==============================================================================
\documentclass[11pt,letterpaper]{article}

%------------------------------- METADATA -------------------------------------
\newcommand{\hwcourse}{EE 999 --- Introduction to Control Systems}
\newcommand{\hwcourseshort}{EE 999}  % running head, keep it short
\newcommand{\hwtitle}{Homework 8}
\newcommand{\hwauthor}{Patrick Gordon}
\newcommand{\hwdate}{\today}
\newcommand{\hwdue}{Due Wednesday, November 11, 2026}
%------------------------------------------------------------------------------

\usepackage[T1]{fontenc}
\usepackage[utf8]{inputenc}
\usepackage{textcomp}
\usepackage[letterpaper,margin=1in,headheight=14pt]{geometry}
\usepackage{amsmath,amssymb}
\usepackage{cancel}
\usepackage{array}
\usepackage{graphicx}
\usepackage[american]{circuitikz}
\usepackage{xcolor}
\usepackage{listings}
\usepackage{fancyhdr}
\usepackage{inconsolata}
\usepackage{pdflscape}
\usepackage{float}
\usepackage[hidelinks]{hyperref}

%--- page furniture -----------------------------------------------------------
\pagestyle{fancy}
\fancyhf{}
\lhead{\hwauthor}
\chead{\hwtitle}
\rhead{\hwcourseshort}
\cfoot{\thepage}

%--- listing styles -----------------------------------------------------------
\definecolor{codebg}{HTML}{F7F7F4}
\definecolor{outbg}{HTML}{FCFCFC}
\definecolor{rule}{HTML}{C8C8C0}
\definecolor{kw}{HTML}{1C4FA1}
\definecolor{cmt}{HTML}{6A7A5A}
\definecolor{str}{HTML}{A0522D}

\lstdefinestyle{mcode}{
  language=Octave,
  basicstyle=\ttfamily\footnotesize,
  keywordstyle=\color{kw}\bfseries,
  commentstyle=\color{cmt}\itshape,
  stringstyle=\color{str},
  backgroundcolor=\color{codebg},
  frame=single,
  rulecolor=\color{rule},
  framesep=6pt,
  numbers=left,
  numberstyle=\ttfamily\scriptsize\color{gray},
  numbersep=8pt,
  showstringspaces=false,
  breaklines=true,
  columns=fullflexible,
  keepspaces=true,
  upquote=true,
}

\lstdefinestyle{moutput}{
  basicstyle=\ttfamily\footnotesize,
  backgroundcolor=\color{outbg},
  frame=single,
  rulecolor=\color{rule},
  framesep=6pt,
  showstringspaces=false,
  breaklines=true,
  columns=fullflexible,
  keepspaces=true,
  upquote=true,
}

%--- structure helpers --------------------------------------------------------
% \problem{1.1b}{Optional title}
\newcommand{\problem}[2]{%
  \bigskip
  \noindent{\large\bfseries Problem #1}%
  \ifx\relax#2\relax\else\ {\large --- #2}\fi
  \par\nobreak\vspace{2pt}\hrule\vspace{8pt}\nobreak
}

% \caplabel{Code}{hw_02_A.m}
\newcommand{\caplabel}[2]{%
  \par\nobreak\vspace{6pt}\nobreak
  \noindent{\small\bfseries #1}\ {\small\ttfamily \detokenize{#2}}%
  \par\nobreak\vspace{2pt}\nobreak
}

% \matlabcode{hw_02_A.m}  -- source listing pulled straight from the .m file
\newcommand{\matlabcode}[1]{%
  \caplabel{Code:}{#1}%
  \lstinputlisting[style=mcode]{#1}%
}

% \codeoutput{output/hw_02_A.txt} -- console output captured by the Makefile
\newcommand{\codeoutput}[1]{%
  \caplabel{Output:}{#1}%
  \IfFileExists{#1}{\lstinputlisting[style=moutput]{#1}}{\missing{#1}{run \texttt{make} to generate it}}%
}

% \answerfigure[width]{figures/foo.pdf}{Caption}  -- placeholder if not exported
% yet. Any format graphicx reads works: pdf, png, jpg.
% A script that plots is exported by the Makefile to figures/<script>.pdf, so
% \answerfigure{figures/hw_02_F.pdf}{...} is how a plotted answer is included.
% The optional argument is the printed width, default the full text width.
% Give a scan or screenshot a smaller one so it is not blown up past its own
% resolution:  \answerfigure[0.8\linewidth]{figures/some_scan.png}{...}
\newcommand{\answerfigure}[3][\linewidth]{%
  \begin{figure}[H]
    \centering
    \IfFileExists{#2}{\includegraphics[width=#1]{#2}}{\missing{#2}{export it, see README}}
    \caption{#3}
  \end{figure}%
}

% \splane{G_1(s)}  -- an empty s-plane grid to sketch poles (x) and zeros (o)
% on by hand. One unit per square, Re from -10 to 3 and Im from -4 to 4.
% \splane[5]{G_2(s)} instead centres the origin, both axes running -5 to 5.
\newcommand{\splane}[2][]{%
  \ifx\relax#1\relax
    \def\spxmin{-10}\def\spxmax{3}\def\spymax{4}\def\spxlabels{-10,-8,...,2}\def\spylabels{-4,-2,2,4}\def\spscale{1}%
  \else
    \def\spxmin{-#1}\def\spxmax{#1}\def\spymax{#1}\def\spxlabels{-#1,...,-1,1,2,...,#1}\def\spylabels{-#1,...,-1,1,2,...,#1}\def\spscale{0.85}%
  \fi
  \begin{figure}[H]
    \centering
    \begin{tikzpicture}[scale=\spscale]
      \draw[gray!30, thin] (\spxmin,-\spymax) grid (\spxmax,\spymax);
      \draw[->, thick] (\spxmin-0.3,0) -- (\spxmax+0.4,0) node[right] {$\mathrm{Re}$};
      \draw[->, thick] (0,-\spymax-0.3) -- (0,\spymax+0.4) node[above] {$\mathrm{Im}$};
      \foreach \x in \spxlabels
        \node[below, font=\scriptsize, fill=white, inner sep=1pt] at (\x,-0.05) {$\x$};
      \foreach \y in \spylabels
        \node[left, font=\scriptsize, fill=white, inner sep=1pt] at (-0.05,\y) {$\y j$};
    \end{tikzpicture}
    \caption{Pole-zero sketch of $#2$.}
  \end{figure}%
}

% \splanewindow{xmin}{xmax}{ymax}{xstep}{ystep}{scale}{Caption}  -- \splane
% with the window chosen per problem, for poles that fall outside the default.
% Re runs xmin..xmax, Im runs -ymax..ymax, axis labels every xstep / ystep, and
% a fine half-unit grid under the unit grid. scale is cm per unit, and the full
% caption is given since the window is usually showing something specific.
%   \splanewindow{-7}{1}{12}{1}{2}{0.6}{Closed-loop poles as $k_d$ varies.}
\newcommand{\splanewindow}[7]{%
  \begin{figure}[H]
    \centering
    \begin{tikzpicture}[scale=#6]
      \draw[gray!15, very thin, step=0.5] (#1,-#3) grid (#2,#3);
      \draw[gray!40, thin] (#1,-#3) grid (#2,#3);
      \draw[->, thick] (#1-0.3,0) -- (#2+0.4,0) node[right] {$\mathrm{Re}$};
      \draw[->, thick] (0,-#3-0.3) -- (0,#3+0.4) node[above] {$\mathrm{Im}$};
      \pgfmathtruncatemacro{\spxa}{#1}\pgfmathtruncatemacro{\spxb}{#1+#4}%
      \foreach \x in {\spxa,\spxb,...,#2}
        {\ifnum\x=0\else\node[below, font=\scriptsize, fill=white, inner sep=1pt] at (\x,-0.05) {$\x$};\fi}
      \pgfmathtruncatemacro{\spya}{-#3}\pgfmathtruncatemacro{\spyb}{-#3+#5}%
      \foreach \y in {\spya,\spyb,...,#3}
        {\ifnum\y=0\else\node[left, font=\scriptsize, fill=white, inner sep=1pt] at (-0.05,\y) {$\y j$};\fi}
    \end{tikzpicture}
    \caption{#7}
  \end{figure}%
}

%--- tables -------------------------------------------------------------------
% Centred paragraph column, so a description column wraps instead of running
% off the page:  \begin{tabular}{|c|P{0.5\linewidth}|}
\newcolumntype{P}[1]{>{\centering\arraybackslash}p{#1}}

% \begin{symboltable}{Signal}{Caption}  -- the symbol / units / description
% table used to define the variables of a problem. The first argument is the
% heading of the left column (Signal, Parameter, State, ...); leave the second
% empty for no caption.
%
%   \begin{symboltable}{Signal}{}
%     \sym{x(t)}{m}{Position of tire center}
%     \sym{r(t)}{m}{Road height (depends on $v_{car}(t)$)}
%   \end{symboltable}
%
% \sym puts its first argument in math mode and brackets the units, so write
% \sym{x(t)}{m}{...}, not \sym{$x(t)$}{[m]}{...}. Anything already math -- a
% subscript, a fraction -- goes in the description as usual: $\frac{1}{4}$.
%
% The optional argument is the width of the description column, for a table
% whose descriptions are much longer or much shorter than average:
%   \begin{symboltable}[0.35\linewidth]{Parameter}{}
\newcommand{\sym}[3]{$#1$ & [#2] & #3 \\ \hline}
\newenvironment{symboltable}[3][0.45\linewidth]{%
  \def\symtabcap{#3}%
  \begin{table}[H]
  \centering
  \renewcommand{\arraystretch}{1.4}
  \begin{tabular}{|c|c|P{#1}|}
  \hline
  #2 & Units & Description \\ \hline\hline
}{%
  \end{tabular}
  \ifx\symtabcap\empty\else\caption{\symtabcap}\fi
  \end{table}%
}

% \begin{answertable}{|c|c|c|}{Caption} -- any other table. Same placement and
% captioning as \answerfigure, but you write the rows yourself:
%
%   \begin{answertable}{|c|P{0.5\linewidth}|}{Pole locations.}
%     \hline
%     Pole & Behaviour \\ \hline\hline
%     $s=-2$ & decays in half a second \\ \hline
%   \end{answertable}
\newenvironment{answertable}[2]{%
  \def\anstabcap{#2}%
  \begin{table}[H]
  \centering
  \renewcommand{\arraystretch}{1.4}
  \begin{tabular}{#1}
}{%
  \end{tabular}
  \ifx\anstabcap\empty\else\caption{\anstabcap}\fi
  \end{table}%
}

% \begin{termtable}{4}{Caption} -- collect like powers of $s$ when multiplying
% out a characteristic polynomial (or any other expansion). The argument is the
% highest power; the columns and the $s^4 \dots s^0$ header are generated from
% it, so changing the degree is a one-character edit. Leave the second argument
% empty for no caption.
%
% Cells are already in math mode: write m_1m_2, not $m_1m_2$. One row per factor
% you multiplied out, one column per power; leave a cell empty where that factor
% contributes nothing to that power. Rows end with \\, like any tabular.
%
%   \begin{termtable}{4}{Terms of the characteristic polynomial.}
%     m_1m_2 & m_1b & m_1k_s       & bk_s       & k_s(k_s+k_w) \\
%            & m_2b & m_2(k_s+k_w) & b(k_s+k_w) &              \\
%            &      & -b^2         & -bk_s      & -k_s^2
%   \end{termtable}
%
% M is a centred column that is always in math mode. It uses \( \) rather than
% $ $ so that an empty cell stays empty -- a bare $$ would start display math.
\newcolumntype{M}{>{\(}c<{\)}}
\makeatletter
\newcount\tt@i
\newcommand{\tt@cols}{}
\newcommand{\tt@head}{}
\newcommand{\tt@build}[1]{%
  \def\tt@cols{}\def\tt@head{}%
  \tt@i=#1\relax
  \@whilenum\tt@i>\m@ne\do{%
    \edef\tt@cols{\tt@cols M}%
    \ifx\tt@head\@empty
      \edef\tt@head{s^{\the\tt@i}}%
    \else
      \edef\tt@head{\tt@head & s^{\the\tt@i}}%
    \fi
    \advance\tt@i\m@ne
  }%
}
\newenvironment{termtable}[2]{%
  \def\tt@cap{#2}%
  \tt@build{#1}%
  \begin{table}[H]%
  \centering
  \setlength{\tabcolsep}{12pt}%
  \renewcommand{\arraystretch}{1.3}%
  \edef\@tempa{\noexpand\begin{tabular}{\tt@cols}}\@tempa
  \tt@head \\ \hline
}{%
  \end{tabular}
  \ifx\tt@cap\@empty\else\caption{\tt@cap}\fi
  \end{table}%
}
\makeatother

% wideproblem environment -- a problem on its own landscape page, centred.
% Use for diagrams much wider than they are tall, where a portrait page would
% shrink the labels below readable size:
%   \begin{wideproblem}{1.1b}{Title}
%     \answerfigure{figures/x.pdf}{Caption}
%   \end{wideproblem}
\newenvironment{wideproblem}[2]{%
  \begin{landscape}%
    \problem{#1}{#2}%
    \vspace*{\fill}%
}{%
    \vspace*{\fill}%
  \end{landscape}%
}

% \prompt{What special property does b*d have?} -- restate the question
\newcommand{\prompt}[1]{%
  \noindent\textit{#1}\par\vspace{6pt}%
}

% \subpart{D.1} -- bold label for a numbered sub-part of a problem
\newcommand{\subpart}[1]{%
  \noindent\textbf{#1}\par\vspace{4pt}%
}

% \answer{...}  -- written answer. Leave the body empty for a visible
% placeholder box; fill it in and the box is replaced by the prose.
\newcommand{\answer}[1]{%
  \ifx\relax#1\relax\todo{write the answer here}\else#1\fi
  \par\vspace{4pt}%
}

\newcommand{\todo}[1]{%
  \fbox{\parbox{0.9\linewidth}{\centering\vspace{10pt}\rmfamily\itshape\small
    #1\vspace{10pt}}}%
}

\newcommand{\missing}[2]{%
  \fbox{\parbox{0.9\linewidth}{\centering\vspace{10pt}\ttfamily\small
    missing: \detokenize{#1}\\[4pt]\rmfamily\itshape #2\vspace{10pt}}}%
}


%==============================================================================
\begin{document}

\begin{center}
  {\LARGE\bfseries \hwtitle}\\[4pt]
  {\large \hwcourse}\\[8pt]
  {\hwauthor} \quad\textbullet\quad {\hwdate}\\[4pt]
  {\small\itshape \hwdue}
\end{center}
\vspace{6pt}\hrule

%------------------------------------------------------------------------------
%  SOLUTIONS
%------------------------------------------------------------------------------


\clearpage
\problem{A}{PI and PD Control}
\answerfigure{figures/fig_hw8_1.png}{A system with PI-type control.\label{fig:hw8-1}}

\problem{A.1}{PI Control}
\prompt{Use $G_c(s) = \dfrac{k_p s + k_i}{s}$ with $k_p = 10$, $k_i = 200$.}

\problem{A.1.a}{}
\prompt{Determine the transfer functions:}

\subpart{i) $T_{ry}(s)$, the input-output transfer function}
\answer{
  \begin{align}
    G_p(s) &= \frac{1/3}{s^2 + 8s + 24} \label{eq:A-Gp} \\
    G_c(s) &= \frac{k_p s + k_i}{s} \label{eq:A-Gc} \\
    Y(s) &= \frac{G_p G_c}{1 + G_p G_c} R(s) \label{eq:A-Y-general} \\
    Y(s) &= \frac{\dfrac{1/3}{s^2 + 8s + 24} \cdot \dfrac{10s + 200}{s}}
                 {1 + \dfrac{1/3}{s^2 + 8s + 24} \cdot \dfrac{10s + 200}{s}} R(s) \label{eq:A-Y-sub} \\
    Y(s) &= \frac{\frac{1}{3}(10s + 200)}
                 {s(s^2 + 8s + 24) + \frac{1}{3}(10s + 200)} R(s) \label{eq:A-Y-cleared} \\
    Y(s) &= \frac{10s + 200}
                 {3s(s^2 + 8s + 24) + 10s + 200} R(s) \label{eq:A-Y-times3} \\
    Y(s) &= \frac{10s + 200}
                 {3s^3 + 24s^2 + 82s + 200} R(s) \label{eq:A-Y-expanded} \\
    T_{ry}(s) = \frac{Y(s)}{R(s)} &= \frac{10s + 200}
                 {3s^3 + 24s^2 + 82s + 200} \label{eq:A-Try}
  \end{align}
}

\subpart{ii) $T_{re}(s)$, the error transfer function}
\answer{
  \begin{align}
    E(s) &= R(s) - Y(s) \label{eq:A-E} \\
    E(s) &= R(s) - T_{ry}(s) R(s) \label{eq:A-E-sub} \\
    E(s) &= \left(1 - \frac{10s + 200}{3s^3 + 24s^2 + 82s + 200}\right) R(s) \label{eq:A-E-factored} \\
    E(s) &= \left(\frac{3s^3 + 24s^2 + 82s + 200}{3s^3 + 24s^2 + 82s + 200} - \frac{10s + 200}{3s^3 + 24s^2 + 82s + 200}\right) R(s) \label{eq:A-E-common} \\
    E(s) &= \frac{3s^3 + 24s^2 + 72s}{3s^3 + 24s^2 + 82s + 200} R(s) \label{eq:A-E-subtracted} \\
    T_{re}(s) = \frac{E(s)}{R(s)} &= \frac{3s^3 + 24s^2 + 72s}{3s^3 + 24s^2 + 82s + 200} \label{eq:A-Tre}
  \end{align}
}

\subpart{iii) $T_{dy}(s)$, the disturbance transfer function}
\answer{
  \begin{align}
    Y(s) &= G_p \left(D(s) + U(s)\right) \label{eq:A-Yd} \\
    Y(s) &= G_p \left(D(s) + G_c(s) E(s)\right) \label{eq:A-Yd-sub} \\
    Y(s) &= G_p \left(D(s) - G_c(s) Y(s)\right) \label{eq:A-Yd-E} \\
    Y(s) &= G_p D(s) - G_p G_c(s) Y(s) \label{eq:A-Yd-dist} \\
    Y(s) + G_p G_c(s) Y(s) &= G_p D(s) \label{eq:A-Yd-collect} \\
    Y(s) \left(1 + G_p G_c(s)\right) &= G_p D(s) \label{eq:A-Yd-factor} \\
    Y(s) &= \frac{G_p}{1 + G_p G_c(s)} D(s) \label{eq:A-Yd-solved} \\
    Y(s) &= \frac{\dfrac{1/3}{s^2 + 8s + 24}}
                 {1 + \dfrac{1/3}{s^2 + 8s + 24} \cdot \dfrac{10s + 200}{s}} D(s) \label{eq:A-Yd-sub-GpGc} \\
    Y(s) &= \frac{\frac{1}{3}s}
                 {s(s^2 + 8s + 24) + \frac{1}{3}(10s + 200)} D(s) \label{eq:A-Yd-cleared} \\
    Y(s) &= \frac{s}
                 {3s(s^2 + 8s + 24) + 10s + 200} D(s) \label{eq:A-Yd-times3} \\
    Y(s) &= \frac{s}
                 {3s^3 + 24s^2 + 72s + 10s + 200} D(s) \label{eq:A-Yd-distributed} \\
    Y(s) &= \frac{s}
                 {3s^3 + 24s^2 + 82s + 200} D(s) \label{eq:A-Yd-expanded} \\
    T_{dy}(s) = \frac{Y(s)}{D(s)} &= \frac{s}
                 {3s^3 + 24s^2 + 82s + 200} \label{eq:A-Tdy}
  \end{align}
}

\problem{A.1.b}{}
\prompt{Draw a pole-zero diagram of $T_{ry}(s)$, label the dominant mode.}
\answer{
\matlabcode{hw_08_A1b.m}
\codeoutput{output/hw_08_A1b.txt}
\answerfigure{figures/hw_08_A1b.pdf}{Pole-zero diagram of $T_{ry}(s)$.}

The dominant poles are at $s = -1.3952 \pm 3.2940j$.
}

\clearpage
\problem{A.2}{}
\prompt{Based on the dominant mode pole locations, determine}
\begin{align}
  p_{1,2} &= -\alpha \pm j\omega \label{eq:A-p12} \\
          &= -1.3952 \pm 3.2940j \label{eq:A-p12-val} \\[\baselineskip]
  \omega_n &= \sqrt{\alpha^2 + \omega^2} \label{eq:A-wn} \\
           &= \sqrt{1.3952^2 + 3.2940^2} \label{eq:A-wn-sub} \\
           &= 3.5773\ \text{rad/s} \label{eq:A-wn-val}
\end{align}

\subpart{i) $t_p$, the peak time}
\answer{
  \begin{align}
    t_p &= \frac{\pi}{\omega} \label{eq:A-tp} \\
        &= \frac{\pi}{3.2940} \label{eq:A-tp-sub} \\
        &= 0.9537\ \text{s} \label{eq:A-tp-val}
  \end{align}
}

\subpart{ii) $\rho$, the damping}
\answer{
  \begin{align}
    \rho &= \frac{\alpha}{\omega_n} \label{eq:A-rho} \\
         &= \frac{1.3952}{3.5773} \label{eq:A-rho-sub} \\
         &= 0.3900 \label{eq:A-rho-val}
  \end{align}
}

\subpart{iii) $M_p$, the percent overshoot}
\answer{
  \begin{align}
    M_p &= 100\% \times e^{-\pi \alpha / \omega} \label{eq:A-Mp} \\
        &= 100\% \times e^{-\pi (1.3952) / 3.2940} \label{eq:A-Mp-sub} \\
        &= 26.43\% \label{eq:A-Mp-val}
  \end{align}
}

\subpart{iv) $t_r$, the rise time}
\answer{
  \begin{align}
    t_r &\approx \frac{1.8}{\omega_n} \label{eq:A-tr} \\
        &\approx \frac{1.8}{3.5773} \label{eq:A-tr-sub} \\
        &\approx 0.5032\ \text{s} \label{eq:A-tr-val}
  \end{align}
}

\subpart{v) $t_s$, the settling time}
\answer{
  \begin{align}
    t_s &\approx \frac{4}{\alpha} \label{eq:A-ts} \\
        &\approx \frac{4}{1.3952} \label{eq:A-ts-sub} \\
        &\approx 2.8670\ \text{s} \label{eq:A-ts-val}
  \end{align}
}

\problem{A.3}{}
\prompt{Analytically determine $e_{ss}$, the steady state error, by determining the system type and $K_{Bode}$.}
\answer{
  \begin{align}
    G_L(s) &= G_c(s)\, G_p(s)\, H_y(s) \label{eq:A-GL} \\
           &= \frac{10s + 200}{s} \cdot \frac{1/3}{s^2 + 8s + 24} \cdot 1 \label{eq:A-GL-sub} \\
           &= \frac{\frac{1}{3}(10s + 200)}{s(s^2 + 8s + 24)} \label{eq:A-GL-mult}
  \end{align}
  $G_L(s)$ has one pole at the origin, so the system is Type~I. A Type~I
  system has no steady-state error for a step input, $e_{ss} = 0$.
}

\clearpage
\problem{A.4}{}
\prompt{Setup $T_{ry}(s)$ in Matlab, run and plot the (true) step response. Determine these values for the complete system}
\begin{align*}
  T_{ry}(s) &= \frac{10s + 200}{3s^3 + 24s^2 + 82s + 200} \tag{\ref{eq:A-Try}}
\end{align*}
\matlabcode{hw_08_A4.m}
\codeoutput{output/hw_08_A4.txt}
\answerfigure{figures/hw_08_A4.pdf}{Step response of $T_{ry}(s)$.}

\subpart{i) $t_p$, the peak time}
\answer{
  \begin{align}
    t_p &= 1.1429\ \text{s} \label{eq:A4-tp}
  \end{align}
}

\subpart{ii) $\rho$, the damping}
\answer{
  \begin{align}
    \rho &= 0.3900 \label{eq:A4-rho}
  \end{align}
}

\subpart{iii) $M_p$, the percent overshoot}
\answer{
  \begin{align}
    M_p &= 20.38\% \label{eq:A4-Mp}
  \end{align}
}

\subpart{iv) $t_r$, the rise time}
\answer{
  \begin{align}
    t_r &= 0.4936\ \text{s} \label{eq:A4-tr}
  \end{align}
}

\subpart{v) $t_s$, the settling time}
\answer{
  \begin{align}
    t_s &= 2.4754\ \text{s} \label{eq:A4-ts}
  \end{align}
}

\subpart{vi) $e_{ss}$, the steady state error}
\answer{
  \begin{align}
    e_{ss} &= 0 \label{eq:A4-ess}
  \end{align}
}

\clearpage
\problem{A.5}{}
\prompt{(Essay question) What is the value of the methods of steps A.1.b--A.2, relative to the method of step A.4?}
\answer{
  The value of the methods in A.1.b--A.2 is that they provide reasonable
  estimates just by evaluating the dominant pole pair without the rigor of
  actually having to plot or otherwise analyze the entire system response.
  The tradeoff is that the numbers are not accurate.
}

\problem{A.6}{}
\prompt{Based on parts A.1.b--A.4, which measures based on the dominant modes came out most accurate? Which measures based on the dominant modes (part 3) came out least accurate?}
\answer{
  \begin{answertable}{|l|c|c|c|}{Dominant-mode estimates (A.2, A.3) compared
    with the complete-system values (A.4). Error is
    $(\text{estimate} - \text{actual}) / \text{actual}$.}
    \hline
    Measure & Estimate & Actual & Error \\ \hline\hline
    $t_p$, peak time        & 0.9537 s & 1.1429 s & $-16.6\%$ \\ \hline
    $\rho$, damping         & 0.3900   & 0.3900   & $0.0\%$ \\ \hline
    $M_p$, percent overshoot & 26.43\%  & 20.38\%  & $+29.7\%$ \\ \hline
    $t_r$, rise time        & 0.5032 s & 0.4936 s & $+1.9\%$ \\ \hline
    $t_s$, settling time    & 2.8670 s & 2.4754 s & $+15.8\%$ \\ \hline
    $e_{ss}$, steady-state error & 0 & 0 & -- \\ \hline
  \end{answertable}

  The damping and steady-state error were the most accurate, and the percent
  overshoot was the least accurate.
}

\clearpage
\problem{B}{PD Control}
\prompt{Use $G_c(s) = k_d s + k_p$ with $k_p = 300$, and $k_d$ to be determined.}
% Same figure as Problem A, keeping its number: the caption steps the counter
% back to where it was, and \thefigure prints the original's number.
{\addtocounter{figure}{-1}\renewcommand{\thefigure}{\ref{fig:hw8-1}}%
\answerfigure{figures/fig_hw8_1.png}{The system from Problem A, repeated, with
  PD control $G_c(s) = k_d s + k_p$ in place of PI.}}

\problem{B.1}{Experiment with $k_d$}

\problem{B.1.a}{}
\prompt{Derive $T_{ry}(s)$, the closed loop transfer function with PD control, and with $k_p$ and $k_d$ as variables. Set $k_p = 300$ and determine the pole locations for several values of $k_d$ on the range from 0.5 to 10.}
\answer{
  \begin{align}
    G_c(s) &= k_d s + k_p \label{eq:B-Gc} \\
    G_p(s) &= \frac{1/3}{s^2 + 8s + 24} \label{eq:B-Gp} \\
    Y(s)   &= \frac{G_p G_c}{1 + G_p G_c} R(s) \label{eq:B-Y-general} \\
    Y(s)   &= \frac{\dfrac{1/3}{s^2 + 8s + 24} \cdot (k_d s + k_p)}
                 {1 + \dfrac{1/3}{s^2 + 8s + 24} \cdot (k_d s + k_p)} R(s) \label{eq:B-Y-sub} \\
    Y(s)   &= \frac{\frac{1}{3}(k_d s + k_p)}
                 {s^2 + 8s + 24 + \frac{1}{3}(k_d s + k_p)} R(s) \label{eq:B-Y-cleared} \\
    Y(s)   &= \frac{k_d s + k_p}
                 {3(s^2 + 8s + 24) + k_d s + k_p} R(s) \label{eq:B-Y-times3} \\
    Y(s)   &= \frac{k_d s + k_p}
                 {3s^2 + 24s + 72 + k_d s + k_p} R(s) \label{eq:B-Y-expanded} \\
    Y(s)   &= \frac{k_d s + k_p}
                 {3s^2 + (24 + k_d)s + (72 + k_p)} R(s) \label{eq:B-Y-collected} \\
    T_{ry}(s) = \frac{Y(s)}{R(s)} &= \frac{k_d s + k_p}
                 {3s^2 + (24 + k_d)s + (72 + k_p)} \label{eq:B-Try}
  \end{align}
\matlabcode{hw_08_B1a.m}
\codeoutput{output/hw_08_B1a.txt}
}

\problem{B.1.b}{}
\prompt{Hand sketch a plot of the closed-loop pole locations as a function of $k_d$.}
\answer{
\splanewindow{-7}{1}{12}{1}{2}{0.6}{Closed-loop poles of $T_{ry}(s)$ as $k_d$ varies
  from 0.5 to 10, with $k_p = 300$.}
}

\problem{B.1.c}{}
\prompt{Select a value of $k_d$ to give $\rho = 0.4$, be rather accurate, e.g., $0.39 < \rho < 0.41$.}
\answer{
\matlabcode{hw_08_B1c.m}
\codeoutput{output/hw_08_B1c.txt}
}

\clearpage
\noindent Using the value of $k_d$ chosen to give $\rho = 0.4$:

\problem{B.2}{}
\prompt{What is the dominant mode of this system? (Same as asking ``where are the dominant pole(s)'')}
\answer{
  With $k_d = 2.7$:
  \begin{align}
    p_{1,2} &= -\alpha \pm j\omega \label{eq:B-p12} \\
            &= -4.4500 \pm 10.2077j \label{eq:B-p12-val}
  \end{align}
}

\problem{B.3}{}
\prompt{Analyze your system (analytically, based on the dominant mode) to determine:}

\subpart{i) $t_p$, the peak time}
\answer{}

\subpart{ii) $\rho$, the damping}
\answer{}

\subpart{iii) $M_p$, the percent overshoot}
\answer{}

\subpart{iv) $t_r$, the rise time}
\answer{}

\subpart{v) $t_s$, the settling time}
\answer{}

\problem{B.4}{}
\prompt{Analytically determine $e_{ss}$, the steady state error, by determining the system type and $K_{Bode}$.}
\answer{}

\clearpage
\problem{B.5}{}
\prompt{Analyze the step response of your system (step response, full system, you can use the \texttt{step} and \texttt{stepinfo} commands) to determine:}

\subpart{i) $t_p$, the peak time}
\answer{}

\subpart{ii) $\rho$, the damping}
\answer{}

\subpart{iii) $M_p$, the percent overshoot}
\answer{}

\subpart{iv) $t_r$, the rise time}
\answer{}

\subpart{v) $t_s$, the settling time}
\answer{}

\subpart{vi) $e_{ss}$, the steady state error (not indicated by \texttt{stepinfo()})}
\answer{}

\problem{B.6}{}
\prompt{Which of the values based on dominant mode are closest to their analysis values? Which are farthest?}
\answer{}

\clearpage
\problem{C}{Compare the PI and PD controllers}
\prompt{Consider these parameters; the desirable characteristic is indicated.}
\begin{answertable}{|l|c|}{}
  \hline
  Controller performance measure & Desirable characteristic \\ \hline\hline
  i) $t_p$, the peak time & Fast \\ \hline
  ii) $\rho$, the damping & Near one \\ \hline
  iii) $M_p$, the percent overshoot & Near zero \\ \hline
  iv) $t_r$, the rise time & Fast \\ \hline
  v) $t_s$, the settling time & Fast \\ \hline
  vi) $e_{ss}$, the steady-state error & Small \\ \hline
\end{answertable}

\problem{C.1}{}
\prompt{What characteristic is most enhanced by the PI control system?}
\answer{}

\problem{C.2}{}
\prompt{What characteristic is most enhanced by the PD control system (tuned for $\rho = 0.4$)?}
\answer{}


\end{document}
